\documentclass[11pt]{article}

\usepackage[margin=1in]{geometry}
\usepackage{amsmath, amssymb}
\usepackage{xcolor}
\usepackage{listings}
\usepackage{hyperref}
\usepackage{array}
\usepackage{booktabs}

\hypersetup{
    colorlinks=true,
    linkcolor=blue!60!black,
    urlcolor=blue!60!black,
    citecolor=blue!60!black
}

\lstdefinestyle{bugpython}{
  language=Python,
  basicstyle=\ttfamily\small,
  keywordstyle=\color{blue!60!black},
  stringstyle=\color{green!40!black},
  commentstyle=\color{gray!70!black},
  showstringspaces=false,
  breaklines=true,
  frame=single,
  rulecolor=\color{gray!30},
  columns=fullflexible,
  keepspaces=true
}

\title{SymPy Bug Report}
\author{}
\date{}

\begin{document}

\maketitle

\section{Bug 5: Wrong Negative-Real Branch in an Antiderivative}

\subsection{Minimal Reproducer}

\medskip

\begin{lstlisting}[style=bugpython]
from sympy import N, diff, integrate, simplify, sqrt, symbols

x = symbols("x")
integrand = sqrt(x**2 + 1)/x
F = integrate(integrand, x)
residual = diff(F, x) - integrand

print("antiderivative:", F)
print("integrand at -2:", integrand.subs(x, -2))
print("derivative at -2:", diff(F, x).subs(x, -2))
print("residual at -2:", simplify(residual.subs(x, -2)))
print("numeric residual:", N(residual.subs(x, -2), 50))
\end{lstlisting}

\subsection{Observed Behavior}

\medskip

\begin{lstlisting}[style=bugpython]
antiderivative: x/sqrt(1 + x**(-2)) - asinh(1/x) + 1/(x*sqrt(1 + x**(-2)))
integrand at -2: -sqrt(5)/2
derivative at -2: sqrt(5)/2
residual at -2: sqrt(5)
numeric residual: 2.2360679774997896964091736687312762354406183596115
\end{lstlisting}

SymPy returns an expression as an antiderivative of
\[
    \frac{\sqrt{x^2 + 1}}{x}.
\]
However, differentiating that returned expression and evaluating at \(x=-2\)
gives \(\sqrt{5}/2\), while the original integrand evaluates to
\(-\sqrt{5}/2\). The residual is therefore \(\sqrt{5}\), so this is not a
harmless formatting difference or an omitted additive constant.

\subsection{Expected Behavior}

The derivative of an indefinite integral result should equal the integrand on
each interval where the integrand is defined. In particular, for
\[
    f(x)=\frac{\sqrt{x^2+1}}{x},
\]
any antiderivative returned for the real variable \(x\) must differentiate back
to \(f(x)\) on both \((-\infty,0)\) and \((0,\infty)\), up to the usual
additive constants on those intervals. At \(x=-2\), the derivative should be
\(-\sqrt{5}/2\), not \(+\sqrt{5}/2\).

\subsection{Mathematical Explanation}

The sign error comes from replacing a branch-sensitive expression by its
positive-real branch. For real nonzero \(x\),
\[
    \frac{\sqrt{x^2+1}}{x}
    =
    \operatorname{sign}(x)\sqrt{1+x^{-2}}.
\]
Thus
\[
    \frac{\sqrt{x^2+1}}{x}
    =
    \sqrt{1+x^{-2}}
    \quad\text{only when } x>0,
\]
whereas on the negative real axis it equals
\[
    -\sqrt{1+x^{-2}}.
\]
SymPy's returned expression differentiates to the positive branch
\(\sqrt{1+x^{-2}}\). That derivative agrees with the integrand for positive
real \(x\), but has the wrong sign for negative real \(x\). Since an additive
constant disappears under differentiation, no choice of integration constant can
repair a pointwise derivative residual of \(\sqrt{5}\) at \(x=-2\).

\subsection{Source-Level Diagnosis}

The diagnosis identifies the relevant integration path as the Meijer-G
indefinite integration fallback. The public call
\(\operatorname{integrate}(\sqrt{x^2+1}/x,x)\) reaches
\texttt{Integral.doit} and \texttt{Integral.\_eval\_integral} in
\nolinkurl{sympy/integrals/integrals.py}. The Risch, heuristic Risch, and
manual integration paths do not produce an answer. The fallback calls
\texttt{meijerint\_indefinite} near
\nolinkurl{sympy/integrals/integrals.py:1105}. The result is then produced in
\nolinkurl{sympy/integrals/meijerint.py} by the internal helper
\texttt{\_meijerint\_}\allowbreak\texttt{indefinite\_1}. The diagnosis
localizes the branch loss around
\nolinkurl{sympy/integrals/meijerint.py:1736}.

The traced rewrite represents the problem through a Meijer-G expression with
argument \(x^2\) and a prefactor \(1/x\). In the diagnosed code region,
\texttt{\_meijerint\_indefinite\_1} extracts powers from the argument and
prefactor, substitutes into the Meijer-G antiderivative, applies
\texttt{hyperexpand}, and then adds a \texttt{powdenest(..., polar=True)}
simplification:

\begin{lstlisting}[style=bugpython]
a, b = _get_coeff_exp(g.argument, x)
_, c = _get_coeff_exp(po, x)
...
fac_ = fac * C * x**(1 + c) / b
...
r = hyperexpand(r.subs(t, a*x**b), place=place)
res += powdenest(fac_*r, polar=True)
\end{lstlisting}

For this integrand, the argument contributes \(b=2\) from \(x^2\), and the
power factor contributes \(c=-1\) from \(1/x\). The generic
substitution--expansion--power-denesting path returns
\[
    \frac{x}{\sqrt{1+x^{-2}}}
    - \operatorname{asinh}\!\left(\frac{1}{x}\right)
    + \frac{1}{x\sqrt{1+x^{-2}}}.
\]
That expression is branch-correct only where
\(\sqrt{1+x^{-2}}\) denotes the same branch as
\(\sqrt{x^2+1}/x\), namely on the positive real branch. On negative real
inputs, the missing \(\operatorname{sign}(x)\) factor causes the derivative to
be \(+\sqrt{1+x^{-2}}\) instead of \(-\sqrt{1+x^{-2}}\). The diagnosis suggests
that this Meijer-G path should avoid applying the rewrite globally when an even
power substitution such as \(x^2\) is combined with an odd prefactor such as
\(1/x\), or should instead return a branch-aware result, preserve polar
conditions through simplification, or decline to evaluate without suitable
assumptions.

\subsection{Additional Failing Instances}

The independent verification checks the same antiderivative identity at the
negative real points \(-1,-2,-3,-4,-5,-6\). These are not different formulas;
they are pointwise witnesses for the same missing negative-real branch. At
every listed point, SymPy's returned antiderivative differentiates to the
positive branch \(\sqrt{1+n^2}/n\), while the original integrand at \(x=-n\)
is the corresponding negative value.

\begin{center}
\renewcommand{\arraystretch}{1.3}
\begin{tabular}{@{}>{\raggedright\arraybackslash}p{0.40\linewidth}>{\raggedright\arraybackslash}p{0.25\linewidth}>{\raggedright\arraybackslash}p{0.25\linewidth}@{}}
\toprule
\textbf{Input / Expression} & \textbf{SymPy behavior} & \textbf{Expected behavior} \\
\midrule
\(x=-1\) in \(\sqrt{x^2+1}/x\) & derivative \(=\sqrt{2}\) & integrand \(=-\sqrt{2}\) \\
\(x=-2\) in \(\sqrt{x^2+1}/x\) & derivative \(=\sqrt{5}/2\) & integrand \(=-\sqrt{5}/2\) \\
\(x=-3\) in \(\sqrt{x^2+1}/x\) & derivative \(=\sqrt{10}/3\) & integrand \(=-\sqrt{10}/3\) \\
\(x=-4\) in \(\sqrt{x^2+1}/x\) & derivative \(=\sqrt{17}/4\) & integrand \(=-\sqrt{17}/4\) \\
\(x=-5\) in \(\sqrt{x^2+1}/x\) & derivative \(=\sqrt{26}/5\) & integrand \(=-\sqrt{26}/5\) \\
\(x=-6\) in \(\sqrt{x^2+1}/x\) & derivative \(=\sqrt{37}/6\) & integrand \(=-\sqrt{37}/6\) \\
\bottomrule
\end{tabular}
\end{center}

The expected behavior is uniform across this family because, for every positive
integer \(n\),
\[
    \left.\frac{\sqrt{x^2+1}}{x}\right|_{x=-n}
    =
    -\frac{\sqrt{n^2+1}}{n}.
\]
The observed derivative has the opposite sign at each of these negative-real
points.

\subsection{Related Issues Assessment}

The bug-finding harness performed a best-effort deduplication check. The
artifact reports the verdict as a new instance of a known failure mode: branch
sensitivity in indefinite integration with radicals and inverse hyperbolic
rewrites is a broadly known kind of problem, especially around transformations
that effectively replace \(\sqrt{x^2}\) by \(x\) rather than \(|x|\). However,
no public report was found for this exact integrand, returned antiderivative, or
pointwise derivative failure at \(x=-2\). The available evidence therefore
supports treating this as an individually actionable SymPy issue and regression
test.

\subsection{Proposed SymPy Regression Test}

\medskip

\begin{lstlisting}[style=bugpython]
import pytest
from sympy import diff, integrate, simplify, sqrt, symbols


@pytest.mark.parametrize("n", [1, 2, 3, 4, 5, 6])
def test_integrate_sqrt_x_squared_plus_one_over_x_negative_branch(n):
    x = symbols("x")
    integrand = sqrt(x**2 + 1)/x
    antiderivative = integrate(integrand, x)
    assert simplify((diff(antiderivative, x) - integrand).subs(x, -n)) == 0
\end{lstlisting}

\end{document}
